\section{Introduction}
\label{sec:introduction}

Large language models are increasingly deployed as autonomous agents that call
external tools and APIs~\citep{schick2024toolformer, yao2023react}. Several
orchestration frameworks structure this behavior as explicit \emph{workflow
graphs}: LangGraph~\citep{langgraph_docs}, CrewAI~\citep{crewai_docs},
AutoGen~\citep{wu2023autogen}, and Google ADK~\citep{google_adk_docs} each
represent computation steps as nodes and allowable transitions as edges. Safety
enforcement in these systems remains almost entirely a \emph{runtime} concern:
tools such as NeMo Guardrails~\citep{nemo_guardrails} and
LlamaGuard~\citep{llamaguard} filter individual calls at execution time,
detecting a violation only when the offending path is exercised.

Some defects, however, are topological. Consider an email-triage agent whose
developer adds a \texttt{draft\_response} node but forgets to connect it to
\texttt{send}: normal-priority emails silently halt at a dead end. A
content-level runtime guardrail cannot catch this unless that path runs, but a
static analysis of the graph flags the dead end and emits a witness trace.
Classical model checkers~\citep{clarke1999modelchecking, holzmann1997spin,
cimatti2002nusmv} verify such properties in principle, but require
hand-translating the system into Promela or SMV---prohibitive for developers
iterating on workflows. For conventional static analyzers, decades of
deployment experience show that what decides adoption is not detection power but
the rate of genuine, actionable findings~\citep{bessey2010coverity,
johnson2013static, sadowski2018lessons}. The analogous question for agent
workflows---how often these graphs actually contain the defects static checks
target, and whether the checks fire on real defects---has received far less
measurement.

\paragraph{Thesis.}
Answering it requires a prior question: \emph{can we trust the graph?} A static
check is only as credible as the behavioral model an extractor recovers from
source, and on public agent code that extraction is lossy---it manufactures
false flags and hides true defects at once. Our central finding is that
\textbf{extraction error dominates check error} and can reverse both apparent
precision and apparent prevalence. We therefore treat static verification of
agent workflows as a \emph{measurement} problem. The extraction-and-checking
machinery we use, \textsc{AgentProof}, is prior, publicly available
infrastructure; this paper's contribution is not that tool but the
cross-framework fidelity measurement it enables and the limits that measurement
exposes.

\paragraph{Study.}
We use \textsc{AgentProof} as a measurement instrument on 922 extracted
file-level graphs mined from public GitHub repositories across the four
frameworks, and validate a stratified 119-graph sample against
\emph{LLM-reconstructed reference graphs} obtained by an independent,
adversarially cross-checked LLM pass. These are \emph{LLM-assisted exploratory
labels}, not human-validated ground truth; a two-human independent validation
protocol (with samples staged) is required future work, and we bound our claims
accordingly. Four questions organize the study: how faithfully workflow
behavior can be recovered from source, and whether correcting the instrument
changes the answer (RQ1); what fraction of raised flags is genuine (RQ2); what
defect prevalence remains once false negatives are audited, and how distinct
estimands differ (RQ3); and whether the static machinery can do a
defect-independent job---certifying soundness and pruning runtime
monitors---without trusting uncertain extraction (RQ4).

\paragraph{Three estimands, not one prevalence.}
Because a lossy extractor both manufactures false flags and hides true defects,
a single pooled ``prevalence'' number conflates distinct quantities. We separate
three (Section~\ref{sec:realworld}, Table~\ref{tab:estimands}): flag-level
precision (what fraction of raised flags are genuine); \emph{structural-defect}
prevalence, a topological property of the graph (1 of 119 sampled graphs); and
\emph{human-gate policy-violation} prevalence, a normative judgment that a
side-effecting tool is reachable with no human checkpoint (3 of 119). A
heterogeneous composite (4 of 119) is reported only as a secondary,
validation-sample rate; post-stratified to the corpus's framework mix it is
2.15\%. All rates are small but not the ${\approx}0$ a na\"ive reading of flag
counts suggests. Three of the four observed cases occur in application-like
files, but the sample is too small to establish a stratum difference (Fisher
exact $p=0.10$), so we make no claim that defects concentrate where they matter.
Correcting a single conditional-edge extractor bug---not any check
algorithm---lifts LangGraph edge recall from 0.68 to 0.82, which is the sense in
which fidelity, not the checkers, is the binding constraint.

\paragraph{Contributions.}
\begin{itemize}\itemsep2pt
  \item \textbf{A cross-framework extraction-fidelity measurement} (LLM-assisted)
  over four frameworks, with a corrected-instrument ablation that isolates
  extraction error from check error and shows a single extractor fix lifting
  LangGraph edge recall 0.68\,$\to$\,0.82 (Section~\ref{sec:realworld}).
  \item \textbf{A separation of three estimands} the field tends to pool:
  flag-level precision, structural-defect prevalence, and human-gate
  policy-violation prevalence, each with repository-clustered uncertainty and
  post-stratified to the corpus composition
  (Section~\ref{sec:realworld}, Table~\ref{tab:estimands}).
  \item \textbf{A taxonomy of extraction failures} and their downstream effect
  on findings, diagnosing fidelity---not the check algorithms---as the binding
  constraint and positioning the study against prior defect-mining work
  (Section~\ref{sec:realworld}, Table~\ref{tab:related}).
  \item \textbf{A conservative certification gate} that returns
  \texttt{inconclusive} rather than a false \texttt{safe} whenever extraction is
  not certified trace-conservative, plus a path-sensitive graph\,$\times$\,DFA
  product that soundly prunes runtime monitors with a small but nonzero gain
  over alphabet-only pruning (Section~\ref{sec:results}, Table~\ref{tab:pruning}).
\end{itemize}
